% ============================================================
%  Rechner-Kochzettel — Casio fx-991DE X (ClassWiz)
%  Komplexe Zahlen & Winkeleinheiten, fachübergreifend
%  Alle Tastenfolgen gegen die offizielle Bedienungsanleitung
%  fx-87DE X / fx-991DE X verifiziert (22.07.2026).
%  ACHTUNG: gilt für das Modell "X", nicht für das neuere "CW".
% ============================================================
\documentclass[a4paper,10pt]{article}
\usepackage[utf8]{inputenc}
\usepackage[T1]{fontenc}
\usepackage[ngerman]{babel}
\usepackage{amsmath,amssymb,mathtools}
\usepackage{xcolor}
\usepackage{geometry}
\usepackage{array}
\usepackage{microtype}
\usepackage{needspace}
\usepackage{tikz}

\geometry{top=10mm,bottom=14mm,left=10mm,right=10mm,includefoot}
\setlength{\parskip}{0pt}
\setlength{\parindent}{0pt}
\renewcommand{\baselinestretch}{0.9}
\setlength{\emergencystretch}{2em}
\hfuzz=3pt
\setlength{\fboxsep}{0.18em}
\setlength{\fboxrule}{0.3pt}
\pagestyle{empty}

\colorlet{cAkz}{cyan!70!black}

% ---- Bereichskopf ----------------------------------------
\newcommand{\bereich}[2][cAkz]{%
  \needspace{4\baselineskip}%
  \vspace{0.5em}%
  \noindent\colorbox{#1!10}{%
    \makebox[\dimexpr\linewidth-2\fboxsep\relax][l]{%
      \small\bfseries\color{#1!80!black}\strut\enspace #2}}%
  \nopagebreak%
  \vspace{0.1em}\par\noindent}

\newcommand{\bereichende}{%
  \vspace{0.2em}%
  {\color{gray!50}\hrule height 0.3pt}%
  \vspace{0.2em}}

% ---- Wenn->Dann-Tabellenlayout ---------------------------
\newenvironment{rezepttab}{%
  \noindent\footnotesize%
  \setlength{\tabcolsep}{0pt}%
  \renewcommand{\arraystretch}{1.0}%
  \begin{tabular}{@{}p{0.34\linewidth}@{\hspace{0.015\linewidth}}p{0.645\linewidth}@{}}%
}{\end{tabular}\par}
\newcommand{\wenn}[1]{\parbox[t]{\linewidth}{\raggedright\scriptsize\itshape #1}}
\newcommand{\dann}[1]{\parbox[t]{\linewidth}{\raggedright #1}}
\newcommand{\rzhl}{\noalign{\vskip 0.18em}\noalign{{\color{gray!40}\hrule height 0.3pt}}\noalign{\vskip 0.18em}}

% ---- Tastensymbol ----------------------------------------
\newcommand{\taste}[1]{\tikz[baseline=(tk.base)]{%
  \node[draw=gray!70,fill=gray!12,rounded corners=1pt,
        inner xsep=2.2pt,inner ysep=1pt,line width=0.3pt,
        font=\scriptsize\sffamily](tk){#1};}}

% ==========================================================
\begin{document}
\begin{center}
  {\large\bfseries Rechner-Kochzettel — Casio fx-991DE X}\\[0.2em]
  {\footnotesize Komplexe Zahlen \& Winkeleinheiten.\;
   Gilt für das Modell \textbf{X} (ClassWiz) — das neuere
   fx-991DE \textbf{CW} hat andere Menüs.}
\end{center}
\vspace{0.3em}

\bereich{R1 \;Vor der Prüfung einmal einstellen}
\begin{rezepttab}
  \wenn{Rechenmodus} &
  \dann{\taste{MENU} \taste{2} \;=\; \glqq Komplexe Zahlen\grqq\
    \quad\scriptsize ohne diesen Modus gibt es kein $i$} \\ \rzhl
  \wenn{Winkeleinheit} &
  \dann{\taste{SHIFT} \taste{MENU} (Setup) $\to$ \emph{Winkeleinheit} $\to$
    \textbf{Gradmaß (D)} oder \textbf{Bogenmaß (R)}\\[0.15em]
    \scriptsize Anzeige oben im Display: \textbf{D} $=$ Grad,
    \textbf{R} $=$ Bogenmaß, \textbf{G} $=$ Gon.
    \emph{Vor jeder Aufgabe kurz hinschauen} — das ist der häufigste
    stille Fehler.} \\ \rzhl
  \wenn{Ergebnisformat komplexer Zahlen} &
  \dann{\taste{SHIFT} \taste{MENU} $\to$ \emph{Komplexe Zahlen} $\to$
    \textbf{a$+$b$i$} (Standard) oder \textbf{r$\angle\theta$}\\[0.15em]
    \scriptsize Anzeige oben: $i$ bei a$+$b$i$, $\angle$ bei Polarform} \\
\end{rezepttab}
\bereichende

\bereich{R2 \;Tasten (im \glqq Komplexe Zahlen\grqq-Modus)}
\begin{rezepttab}
  \wenn{$i$ eingeben \;— in der E-Technik heißt es $j$, dasselbe} &
  \dann{\taste{ENG} \quad\scriptsize violettes $i$ auf der Taste} \\ \rzhl
  \wenn{$\angle$ für Polareingabe, z.\,B.\ $200\angle 0$} &
  \dann{\taste{SHIFT} \taste{ENG}} \\ \rzhl
  \wenn{\textbf{Realteil} (ReP)} &
  \dann{\taste{OPTN} \taste{3} \quad\scriptsize liefert die nackte Zahl
    ohne $i$, direkt weiterrechenbar} \\ \rzhl
  \wenn{\textbf{Imaginärteil} (ImP)} &
  \dann{\taste{OPTN} \taste{4} \quad\scriptsize ebenfalls ohne $i$} \\ \rzhl
  \wenn{\textbf{Betrag} $|\underline Z|$ (Abs)} &
  \dann{\taste{SHIFT} \taste{(}} \\ \rzhl
  \wenn{\textbf{Argument} / Phasenwinkel (Arg)} &
  \dann{\taste{OPTN} \taste{1}} \\ \rzhl
  \wenn{Konjugierte $\underline Z^{*}$ (Conjg)} &
  \dann{\taste{OPTN} \taste{2}} \\ \rzhl
  \wenn{\emph{ein} Ergebnis in Polarform sehen} &
  \dann{\taste{OPTN} \taste{$\blacktriangledown$} \taste{1}
    \;($\blacktriangleright r\angle\theta$)} \\ \rzhl
  \wenn{\emph{ein} Ergebnis in a$+$b$i$ sehen} &
  \dann{\taste{OPTN} \taste{$\blacktriangledown$} \taste{2}
    \;($\blacktriangleright a+bi$)} \\
\end{rezepttab}
\bereichende

\bereich{R3 \;Typische Eingaben (Wechselstrom- und Leitungsrechnung)}
\begin{rezepttab}
  \wenn{Impedanz eintippen, z.\,B.\ $(200-j174{,}4)\,\Omega$} &
  \dann{\taste{2}\taste{0}\taste{0} \taste{$-$}
    \taste{1}\taste{7}\taste{4}\taste{.}\taste{4} \taste{ENG}
    \taste{$=$}} \\ \rzhl
  \wenn{durch eine komplexe Zahl \textbf{teilen}} &
  \dann{Nenner \textbf{immer klammern}:\;
    $a$ \taste{$\div$} \taste{(} $\dots$ \taste{ENG} \taste{)}
    \taste{$=$}\\[0.15em]
    \scriptsize ohne Klammer teilt der Rechner nur durch den Realteil und
    addiert den Rest — der häufigste Tippfehler überhaupt} \\ \rzhl
  \wenn{Quadrat durch komplexe Zahl,
    z.\,B.\ $\underline Z=\tfrac{Z_L^2}{\underline Z(0)}$
    ($\lambda/4$-Trafo)} &
  \dann{$Z_L$ \taste{$x^2$} \taste{$\div$} \taste{(} $\underline Z(0)$
    \taste{)} \taste{$=$}
    \quad\scriptsize Real- und Imaginärteil direkt ablesen} \\ \rzhl
  \wenn{Kehrwert: Admittanz $\underline Y=\tfrac{1}{\underline Z}$
    (Parallelschaltung)} &
  \dann{\taste{(} $\underline Z$ \taste{)} \taste{$x^{-1}$};\quad
    $\underline Y$-Werte addieren, am Schluss wieder \taste{$x^{-1}$}
    zurück nach $\underline Z$} \\ \rzhl
  \wenn{nur den Real- \emph{oder} Imaginärteil in die nächste
    Rechnung übernehmen} &
  \dann{\taste{OPTN} \taste{4} \taste{Ans} \taste{)} \taste{$=$}
    $\;\to\;$ Imaginärteil als gewöhnliche reelle Zahl, die selbst wieder
    \taste{Ans} ist (Realteil: \taste{OPTN} \taste{3})\\[0.15em]
    \scriptsize Achtung: jedes \taste{$=$} überschreibt \taste{Ans}.
    Brauchst du \emph{beide} Teile, vorher mit \taste{STO} in einer
    Variablen ablegen — oder schlicht beide Zahlen notieren.} \\ \rzhl
  \wenn{$\tan\beta l$ und andere Winkelfunktionen} &
  \dann{Argument in \textbf{rad} eintippen $\Rightarrow$ Setup auf
    \textbf{R}; in Grad eintippen $\Rightarrow$ Setup auf \textbf{D}.
    Nur nie mischen.} \\ \rzhl
  \wenn{dB-Rechnungen: $\lg$ und $\ln$} &
  \dann{\taste{log} ohne Basisangabe $=$ Basis 10 $=\lg$;\quad
    \taste{ln} ist die eigene Taste für $\ln$} \\
\end{rezepttab}
\bereichende

\bereich{R4 \;Fallen, die still danebenrechnen}
\begin{rezepttab}
  \wenn{Winkeleinheit wirkt \textbf{doppelt}} &
  \dann{dieselbe Einstellung gilt für die \emph{Eingabe} in
    $\sin/\cos/\tan$ \textbf{und} für die \emph{Ausgabe} von Winkeln
    (Arg, Polarform). Rechnung in rad, Ergebnis soll in Grad
    $\Rightarrow$ vor dem Ablesen umstellen.} \\ \rzhl
  \wenn{Phasenwinkel sieht \glqq falsch\grqq\ aus} &
  \dann{$\theta$ wird nur im Bereich $-180^\circ<\theta\le180^\circ$
    angezeigt — aus $216^\circ$ wird $-144^\circ$.\\[0.15em]
    \scriptsize Für $\tan$ egal (Periode $180^\circ$), beim abzulesenden
    Phasenwinkel nicht.} \\ \rzhl
  \wenn{$i$ lässt sich nicht eingeben} &
  \dann{falscher Modus — \taste{MENU} \taste{2}. Das violette $i$ ist nur
    im \glqq Komplexe Zahlen\grqq-Modus erreichbar.} \\ \rzhl
  \wenn{nach jedem Moduswechsel} &
  \dann{Blick auf die Indikatoren oben im Display
    (\textbf{D}/\textbf{R}/\textbf{G} und $i$/$\angle$) —
    kostet zwei Sekunden und rettet ganze Teilaufgaben} \\
\end{rezepttab}
\bereichende

\vfill
{\scriptsize\color{gray!60!black}
Tastenfolgen verifiziert gegen die offizielle Bedienungsanleitung
\emph{fx-87DE X / fx-991DE X} (Casio), Stand 22.07.2026.}

\end{document}
